\clearpage
\fancyhead[L]{\small HARDING. HIMMELSKARTEN.}
\hypertarget{native-598-lower}{}
\begin{center}\small Göttingische gelehrte Anzeigen. Stück 5. Seite 41..45. 1817 Januar 9.\\[3pt]\rule{.72\linewidth}{.3pt}\end{center}

Von den Himmelskarten unsers Herrn Prof. \textsc{Harding}, wovon wir die letzten Lieferungen im 158. Stück dieser Blätter von 1812 [S. 541 d. B.] angezeigt haben, ist jetzt wiederum in Commission der Vandenhoeck-Ruprechtschen Buchhandlung eine neue, die \emph{fünfte}, Lieferung von vier Blättern erschienen, nemlich die Nummern VI, X, XIX und XXV. Um leichter zu übersehen, ein wie grosser Theil dieser nützlichen Arbeit jetzt vollendet ist, und wie viel noch fehlt, bringen wir hier nur kürzlich die Veranlassung und den Plan derselben in Erinnerung. Die Entdeckung der Planeten Ceres und Pallas, die nie grösser als wie Fixsterne achter Ordnung, oft noch viel kleiner, erscheinen, machte zuerst weit detaillirtere Sternkarten, als wir bis dahin besassen, von denjenigen Theilen des Himmels, in welchen jene interessanten Himmelskörper sich zeigen können, unentbehrlich. Um nicht alle Jahre für dieselben besonderer Karten zu bedürfen, war ein eigener Atlas für den Zodiakus dieser Planeten wünschenswerth, den der Verfasser zu bearbeiten beschloss. Eine mässige Anzahl von Blättern (10 oder 11) würde hiezu hingereicht haben. Allein während der Verf. diese Arbeit bereits verfolgte, und zum Theil eben durch diese Arbeit veranlasst, kam noch die Entdeckung zweier andern neuen Planeten hinzu, die eine Erweiterung des Planes nothwendig machten. Der mannichfaltige Nutzen, welchen detaillirte Stern\SourceLeaf{599}karten auch in andern Beziehungen verschaffen, bewog den Verfasser, seinen Atlas, von welchem inzwischen schon einige Blätter vollendet waren, über den ganzen nördlichen Himmel, und denjenigen Theil des südlichen, der in unsern Gegenden von Europa noch gut beobachtet werden kann, und wozu noch die Materialien vorhanden sind, auszudehnen, also bis etwas mehr als $30^\circ$ südliche Abweichung. Neun Blätter, jedes von $40^\circ$ Ausdehnung in der geraden Aufsteigung, sollten also diesen südlichen Theil des Himmels, und eben so neun Blätter, jedes von ähnlicher Ausdehnung, den Theil des nördlichen Himmels vom Aequator bis zu $32^\circ$ nördlicher Abweichung umfassen. Diese 18 Blätter sind jetzt, nachdem die bisher noch fehlenden VI und X in vorliegender Lieferung hinzugekommen sind, ganz vollendet. Nur wird in Zukunft das Blatt Nr. I., welches zu Anfang und noch nach dem beschränkten Plane ausgearbeitet, von beinahe $22^\circ$ südlicher bis fast $13^\circ$ nördlicher Abweichung sich erstreckt, und dessen nördlicher Theil auch auf dem neuen Blatt Nr. X. mit enthalten ist, ausgeschlossen und durch ein anderes bis $32^\circ$ südliche Abweichung gehendes ersetzt werden, welches demnächst den Besitzern des Atlasses von dem Verfasser unentgeldlich nachgeliefert werden wird.

Der Raum des Himmels, welchen diese 18 Blätter umfassen, ist für die Astronomen in so fern der interessanteste, weil alle Planeten und der Mond nie über denselben hinausgehen, und desshalb glaubte der Verfasser mit Recht, diesen zuerst vollenden zu müssen. Allein eine eben so detaillirte Bearbeitung des nördlichen Himmels von $32^\circ$ Abweichung bis zum Pole würde schon wegen der so häufig denselben durchstreifenden Cometen als ein wesentliches Bedürfniss angesehen werden müssen, da so manche Cometen \emph{nur} teleskopisch sind, und selbst die grössern immer noch mit Fernröhren eine Zeitlang verfolgt werden können, wenn ihre zunehmende Entfernung sie bereits dem unbewaffneten Auge entzogen hat. Bei solchen Veranlassungen muss der Beobachter in Ermangelung gestochener Sternkarten immer sich der zeitraubenden Arbeit unterziehen, sich selbst von der Gegend, die der Comet durchläuft, eine Karte zu entwerfen. Die Blätter des vorliegenden Atlasses XIX und XXV machen den Anfang, diesem Bedürfnisse abzuhelfen. Der Plan zur Abtheilung des nördlichen Himmels, der hiebei befolgt wird, ist folgender: Da die Grade der Parallelkreise näher nach dem Pole zu immer kleiner werden, so konnte ein Blatt von derselben Grösse, wie die ersten 18, mehr in Rectascension fassen. Acht Blätter, jedes unten von etwas über $45^\circ$ Ausdehnung in gerader Aufsteigung, und in der Mitte von $30^\circ$ bis fast $65^\circ$ Abweichung sich erstreckend, werden zuerst folgen; der noch übrige Theil des Himmels von $64^\circ$ Abweichung bis zum Pole wird dann bei nur wenig verkleinertem Massstabe in zwei Blättern abzuthun sein, auf denen die Ecken, in welchen sonst schon auf den Blättern vorgekommene Theile des Himmels wiederholt werden müssten, zur Darstellung interessanter Sterngruppen, wie der Plejaden, der Krippe und anderer, verwandt werden sollen. In 28 Blättern wird also der ganze Athas vollendet sein, und die jetzt noch fehlenden acht Blätter, die schon grösstentheils gezeichnet sind, werden hoffentlich in kurzer Zeit geliefert werden können.

Bei den ersten 18 Blättern war eine künstlichere Projectionsart überflüssig; das Netz aus blossen Quadraten bestehend erleichterte nicht allein die Zeichnung, sondern gewährt auch bei dem Gebrauche manche Bequemlichkeit, wogegen die kleine ohnehin kaum merkliche Abweichung von dem wahren Verhältniss der Rectascensions- und Declinationsgrade, gegen die Grenze zu, in gar keinen Betracht kam. Allein bei grössern Declinationen war dies nicht mehr anwendbar. Für die acht Blätter XIX bis XXVI hat also der Verf. eine ihm von Hrn. Hofr. \textsc{Gauss} dazu vorgeschlagene Projectionsart gewählt, bei welcher die Declinationskreise gerade Linien, die Parallelkreise concentrische Kreise sind, deren abnehmende Halbmesser nach aller Schärfe so bestimmt sind, dass überall die Rectascensionsgrade zu den Declinationsgraden ihr richtiges Verhältniss bekommen. Bei dieser Darstellung ist also die Abbildung dem Himmel in den kleinsten Theilen vollkommen ähnlich, welches die wesentlichste Bedingung ist, \SourceLeaf{600}Die Declinationsgrade sind freilich genau genommen, nicht durchaus von gleicher Grösse. Allein dieser, der Natur der Sache nach unvermeidliche Umstand ist hier vollkommen gleichgültig, zumal da der Unterschied so gering ist, dass er kaum bemerkt wird: auf dem 64sten und 3osten Grade der Declination beträgt nemlich ein Declinationsgrad 13,07 Millimeter, in der Mitte auf dem 47sten Grade hingegen 12,38 Millimeter.

Da die Vertheilung des Himmels auf die Karten nur durch die geraden Aufsteigungen und Abweichungen bestimmt ist, so erscheinen natürlich manche Sternbilder zerstückelt. Das Blatt Nr.\ VI enthält nemlich theilweise die Jungfrau, Waage, den Scorpion, die Schlange, den Wolf, Centaur, Schwanz der Hydra und den ganzen Vogel Einsiedler; das Blatt Nr.\ X theilweise die Fische, die Andromeda, die Dreiecke, den Widder, Wallfisch und Pegasus; das Blatt Nr.\ XIX theilweise den Cepheus, die Cassiopea, den Erntehüter, Camelopard, Perseus, die Andromeda, die Fische, die Dreiecke und den Widder; das Blatt Nr.\ XXV theilweise den Drachen, den Herkules, die Leyer, den Schwan und den Cepheus. Papier und Stich stehen an Schönheit den frühern Lieferungen nicht nach.
